\begin{scanpage}{0199}{181}
It is not difficult to see that in the case of a passive vibrator the reflection of current waves at its ends occurs in the same way as at the end of a semi-infinite conductor.

Thus, the complex amplitudes of current waves in a thin, finite length conductor are proportional to the functions $\psi(z)$ and $\psi_\pm(z)$ which monotonically decrease with an increase of $z$ as a consequence of radiation. Let us note several properties of current waves in a vibrator. Each advancing wave in sum with the reflected wave excited by it gives a zero current at the conductor's end. In the case $L=z_2-z_1\approx n\frac\lambda2$ $(n=1,2,3,\ldots)$ and $D\approx0$, a current resonance begins in the vibrator.

The precision of Expressions (28.03) and (28.14) obtained by the method of slowly varying functions is different in various sections of the conductor. It is comparatively low near the conductor's ends (and in the vicinity of the point $z=0$ of a transmitting vibrator) where the current waves arise, and where their complex amplitude varies rather rapidly. As the distance from these vibrator elements increases, the precision of these equations increases without limit.

It should be stated that with a more rigorous approach [79, 90] the amplitudes of all the reflected waves will be determined by different functions; however, the difference between them rapidly decreases with an increase of the reflection number. The functions $\psi(z)$ and $\psi_\pm(z)$ only approximately describe the current waves, but on the other hand they allow one to effectively sum them and to obtain closed equations.

Using the variational method for the functions $\psi(z)$ and $\psi_\pm(z)$, we obtained the approximate, but on the other hand, simple equations (see [83])
\[
\psi(z)=\frac{\displaystyle\ln\frac{-1}{\gamma^2q}}{\displaystyle\ln\frac{2ix}{\gamma}-E(2qx)e^{-2iqx}},\qquad\Biggr\}\tag{28.21}
\]
\begin{flushright}(Equation continued on next page)\end{flushright}
\end{scanpage}
